\chapter{Unit 3: Biodiversity Indices, Conservation \& Protected Areas}

\begin{formulabox}{Quantitative Diversity Indices}
\begin{itemize}[leftmargin=1.2em]
    \item \textbf{Species Richness ($S$)}: Total absolute number of distinct species recorded in a defined sample or biotope.
    \item \textbf{Shannon-Wiener Information Diversity Index ($H'$)}:
    \begin{equation}
    H' = -\sum_{i=1}^S p_i \ln(p_i)
    \end{equation}
    where $p_i = n_i / N$ is the relative abundance proportion of species $i$.
    \item \textbf{Maximum Diversity \& Pielou's Evenness ($J'$)}:
    \begin{align}
    H'_{\text{max}} &= \ln(S) \\
    J' &= \frac{H'}{H'_{\text{max}}} = \frac{H'}{\ln(S)}, \quad 0 \le J' \le 1.0
    \end{align}
    \item \textbf{Simpson's Dominance Index ($D$)}:
    \begin{equation}
    D = \sum_{i=1}^S p_i^2 = \frac{\sum_{i=1}^S n_i(n_i - 1)}{N(N - 1)}
    \end{equation}
    \item \textbf{Simpson's Index of Diversity ($1 - D$)} and \textbf{Reciprocal Simpson ($1/D$)}.
    \item \textbf{Alpha, Beta, and Gamma Diversity (Whittaker, 1960)}:
    \begin{align}
    \alpha\text{-diversity} &= \text{Local within-habitat species richness} \\
    \beta\text{-diversity} &= \frac{\gamma}{\bar{\alpha}} = \text{Species turnover rate between habitats across gradients} \\
    \gamma\text{-diversity} &= \text{Total landscape or regional biodiversity}
    \end{align}
\end{itemize}
\end{formulabox}

\begin{formulabox}{Hotspots, Megadiversity \& Biogeographic Realms}
\begin{itemize}[leftmargin=1.2em]
    \item \textbf{Myers Biodiversity Hotspot Criteria (2000)}:
    \begin{enumerate}
        \item Must contain at least $1{,}500$ endemic species of vascular plants ($>0.5\%$ of world total).
        \item Must have lost at least $70\%$ of its primary original native vegetation ($\le 30\%$ remaining intact).
    \end{enumerate}
    \item \textbf{The 4 Indian Biodiversity Hotspots}:
    \begin{enumerate}
        \item \textbf{Western Ghats and Sri Lanka}: Exceptionally high amphibian, reptilian, and floral endemism.
        \item \textbf{Himalaya}: Entire Indian Himalayan mountain arc (spanning Jammu \& Kashmir, HP, Uttarakhand, Sikkim, Arunachal Pradesh).
        \item \textbf{Indo-Burma}: Covers North-Eastern states south of the Brahmaputra River, Andaman Islands.
        \item \textbf{Sundaland}: Includes the Nicobar group of islands (extending through Malaysia and Indonesia).
    \end{enumerate}
    \item \textbf{Megadiverse Countries}: India is recognized among the 17 megadiverse nations, harboring $\approx 7\text{--}8\%$ of global recorded floral and faunal species across only $2.4\%$ of global terrestrial land area.
    \item \textbf{Biogeographic Zones of India (Rodgers \& Panwar Classification - 10 Zones)}:
    1. Trans-Himalaya, 2. Himalaya, 3. Desert, 4. Semi-Arid, 5. Western Ghats, 6. Deccan Peninsula, 7. Gangetic Plain, 8. Coasts, 9. North-East, 10. Islands.
\end{itemize}
\end{formulabox}

\begin{lawbox}{IUCN Categories, Protected Area Frameworks \& Statutes}
\begin{itemize}[leftmargin=1.2em]
    \item \textbf{IUCN Red List Categories}: Extinct (EX), Extinct in the Wild (EW), \textbf{Critically Endangered (CR)}, \textbf{Endangered (EN)}, \textbf{Vulnerable (VU)}, Near Threatened (NT), Least Concern (LC), Data Deficient (DD).
    \item \textbf{Quantitative Thresholds for Threatened Categories}:
    \begin{center}
    \begin{tabularx}{\linewidth}{X c c c}
    \toprule
    \textbf{IUCN Quantitative Criterion} & \textbf{Critically Endangered (CR)} & \textbf{Endangered (EN)} & \textbf{Vulnerable (VU)} \\
    \midrule
    Population reduction over 10 yrs & $\ge 90\%$ (observed) & $\ge 70\%$ & $\ge 50\%$ \\
    Extent of Occurrence (EOO) & $< 100\text{ km}^2$ & $< 5{,}000\text{ km}^2$ & $< 20{,}000\text{ km}^2$ \\
    Area of Occupancy (AOO) & $< 10\text{ km}^2$ & $< 500\text{ km}^2$ & $< 2{,}000\text{ km}^2$ \\
    Number of mature individuals & $< 50$ & $< 250$ & $< 1{,}000$ \\
    \bottomrule
    \end{tabularx}
    \end{center}
    \item \textbf{In-Situ vs Ex-Situ Conservation}:
    \begin{itemize}
        \item \textbf{In-Situ}: National Parks (no human grazing/habitation), Wildlife Sanctuaries (limited rights permitted), Biosphere Reserves (Core, Buffer, Transition zones), Tiger Reserves, Sacred Groves.
        \item \textbf{Ex-Situ}: Zoological gardens, Botanical gardens, Seed banks, Cryopreservation, In-vitro gene banks.
    \end{itemize}
    \item \textbf{Key Wildlife Statutes}: Wildlife (Protection) Act, 1972 (amended 2022); Biological Diversity Act, 2002 (National Biodiversity Authority - NBA, State Biodiversity Boards - SBB, Biodiversity Management Committees - BMC).
\end{itemize}
\end{lawbox}
